% Ad Familiares 8.8 --- headnote
% ad-familiares-08-08, early October 51 BC, Rome

\textit{M. Caelius Rufus to Cicero, written from Rome
in the early part of October 51 BC (the manuscript
dateline: \textit{Scr. Romae in. m. Oct. a. 703 (51)}).
The longest and most substantive of the surviving
Caelius newsletters, and the political pivot of book 8.
The previous dispatches had reported a year of
deadlock --- Marcellus's procrastinations, Pompey's
elusiveness, quorum failures, vetoes signalled but not
yet cast. The first two paragraphs continue in the
familiar Forum-vignette register (the prosecution of
Sempronius Rufus under the calumnia-statute; the
strange aborted extortion-trial of M. Servilius, where
the praetor Laterensis bungled the count and the
defendant ended neither acquitted nor condemned), with
Caelius doing what he does best: dramatising himself
into the centre of every scene.}

\textit{Then, in section 4, the letter turns: the
provinces have at last been debated, Pompey has shown
his hand, and Caelius proudly transcribes for Cicero
the full text of the senatus consultum of the day
before the Kalends of October --- four resolutions
in all, with the witnesses listed at the head and the
tribunes who interposed listed at the foot. This is the
decree that fixes the Ides of March 50 BC as the date
for revisiting Caesar's Gallic command, and it is the
first formal step on the road to civil war; Caelius
seems to know it. The ledger-prose of the consultum
text and the bureaucratic list of attesting senators
contrast deliberately with the Forum-narrative voice
around them: Caelius is showing Cicero his proper
documentary care. Pompey's two reported \textit{vocones}
in section 9 --- that after the Kalends of March he
will not hesitate, and that he sees no difference
between a Caesar who disobeys the Senate and a Caesar
who organises men to prevent the Senate from decreeing
--- are the most explicit public statement on record
that the breach between him and Caesar is now real.
The letter closes on Caelius's standing private
request --- panthers for his aedilician games --- and
the Sittian bond Cicero is asked to look after on his
behalf.}
